% Game rating: the paper. What is typed here says what the numbers mean.
% Every number, table, figure and board is written by `lab paper game-rating`
% into generated/, from the model's committed results, and read back with
% \V{key} and \block{name}.
\documentclass[11pt]{article}
% The preamble every paper of the lab shares: % The preamble every paper of the lab shares: % The preamble every paper of the lab shares: \input{../../../papers/preamble}
% from models/<model>/paper/, \input{../preamble} from papers/<name>/.
% Boards come from chessboard and xskak, plots from pgfplots. A paper prints no
% number by hand: `lab paper` writes every number, table and figure under the
% paper's generated/ folder, and paper.tex says what they mean.
\usepackage[LSBC4,T1]{fontenc}
\usepackage[utf8]{inputenc}
\usepackage{lmodern}
\usepackage[margin=1in]{geometry}
\usepackage{microtype}
\usepackage{amsmath}
\usepackage{booktabs}
\usepackage{tabularx}
\usepackage{array}
\usepackage{xcolor}
\usepackage[obeyspaces,spaces]{url}
\usepackage{xskak}
\usepackage{pgfplots}
\usepgfplotslibrary{groupplots}
\usetikzlibrary{positioning,arrows.meta,calc,fit}
\pgfplotsset{compat=1.18}
\usepackage{caption}
\usepackage{subcaption}
\usepackage{enumitem}
\usepackage[section]{placeins}
\usepackage[hidelinks]{hyperref}

% A generated value by its key: generated/values.tex sets them, \V{key} prints
% one, and a key nobody set stops the build.
\makeatletter
\newcommand{\setvalue}[2]{\expandafter\gdef\csname labvalue@#1\endcsname{#2}}
\newcommand{\V}[1]{\@ifundefined{labvalue@#1}%
  {\PackageError{lab}{No generated value named '#1'}{Run `lab paper` for this model again.}}%
  {\csname labvalue@#1\endcsname}}
\makeatother

% Chart colours. A run keeps its colour in every figure of every paper.
\definecolor{labBlue}{HTML}{2A78D6}
\definecolor{labOrange}{HTML}{EB6834}
\definecolor{labAqua}{HTML}{1BAF7A}
\definecolor{labYellow}{HTML}{EDA100}
\definecolor{labMagenta}{HTML}{E87BA4}
\definecolor{labGreen}{HTML}{008300}
\definecolor{labViolet}{HTML}{6250D6}
\definecolor{labGray}{HTML}{898781}
\definecolor{labInk}{HTML}{0B0B0B}
\definecolor{labSecondary}{HTML}{52514E}
\definecolor{labGrid}{HTML}{E1E0D9}
\definecolor{labAxis}{HTML}{C3C2B7}
\definecolor{labWash}{HTML}{CDE2FB}

\pgfplotsset{
  lab/.style={
    font=\sffamily\footnotesize,
    axis line style={labAxis},
    tick style={labAxis},
    tick label style={labSecondary},
    label style={labSecondary},
    title style={labInk, font=\sffamily\small\bfseries, align=left},
    grid style={labGrid, line width=0.3pt},
    major grid style={labGrid, line width=0.3pt},
    axis lines*=left,
    tick align=outside,
    clip=false,
    legend style={draw=none, fill=none, font=\sffamily\footnotesize, text=labSecondary,
      cells={anchor=west}, legend columns=-1, column sep=6pt, at={(0,1.03)}, anchor=south west},
    every axis plot/.append style={line width=1.1pt, line join=round, line cap=round},
    scaled ticks=false,
    /pgf/number format/1000 sep={,},
  },
  lab marks/.style={mark=*, mark size=1.9pt, mark options={draw=white, line width=0.6pt}},
}

\setchessboard{boardfontencoding=LSBC4, boardfontfamily=skaknew, blackfieldcolor=black!18,
  whitefieldcolor=white, setfontcolors, showmover=false, label=true, labelfontsize=6pt,
  boardfontsize=12pt}

% A move or line in UCI, a position in FEN (breakable at a slash), and a run's name.
\newcommand{\uci}[1]{\texttt{#1}}
\newcommand{\fen}[1]{{\footnotesize\path{#1}}}
\newcommand{\run}[1]{\texttt{#1}}
\newcolumntype{L}{>{\raggedright\arraybackslash}X}
\newcolumntype{R}{>{\raggedleft\arraybackslash}X}

\setlength{\parskip}{0.4em}
\setlength{\parindent}{0pt}
\captionsetup{font=small, labelfont=bf}

% from models/<model>/paper/, \input{../preamble} from papers/<name>/.
% Boards come from chessboard and xskak, plots from pgfplots. A paper prints no
% number by hand: `lab paper` writes every number, table and figure under the
% paper's generated/ folder, and paper.tex says what they mean.
\usepackage[LSBC4,T1]{fontenc}
\usepackage[utf8]{inputenc}
\usepackage{lmodern}
\usepackage[margin=1in]{geometry}
\usepackage{microtype}
\usepackage{amsmath}
\usepackage{booktabs}
\usepackage{tabularx}
\usepackage{array}
\usepackage{xcolor}
\usepackage[obeyspaces,spaces]{url}
\usepackage{xskak}
\usepackage{pgfplots}
\usepgfplotslibrary{groupplots}
\usetikzlibrary{positioning,arrows.meta,calc,fit}
\pgfplotsset{compat=1.18}
\usepackage{caption}
\usepackage{subcaption}
\usepackage{enumitem}
\usepackage[section]{placeins}
\usepackage[hidelinks]{hyperref}

% A generated value by its key: generated/values.tex sets them, \V{key} prints
% one, and a key nobody set stops the build.
\makeatletter
\newcommand{\setvalue}[2]{\expandafter\gdef\csname labvalue@#1\endcsname{#2}}
\newcommand{\V}[1]{\@ifundefined{labvalue@#1}%
  {\PackageError{lab}{No generated value named '#1'}{Run `lab paper` for this model again.}}%
  {\csname labvalue@#1\endcsname}}
\makeatother

% Chart colours. A run keeps its colour in every figure of every paper.
\definecolor{labBlue}{HTML}{2A78D6}
\definecolor{labOrange}{HTML}{EB6834}
\definecolor{labAqua}{HTML}{1BAF7A}
\definecolor{labYellow}{HTML}{EDA100}
\definecolor{labMagenta}{HTML}{E87BA4}
\definecolor{labGreen}{HTML}{008300}
\definecolor{labViolet}{HTML}{6250D6}
\definecolor{labGray}{HTML}{898781}
\definecolor{labInk}{HTML}{0B0B0B}
\definecolor{labSecondary}{HTML}{52514E}
\definecolor{labGrid}{HTML}{E1E0D9}
\definecolor{labAxis}{HTML}{C3C2B7}
\definecolor{labWash}{HTML}{CDE2FB}

\pgfplotsset{
  lab/.style={
    font=\sffamily\footnotesize,
    axis line style={labAxis},
    tick style={labAxis},
    tick label style={labSecondary},
    label style={labSecondary},
    title style={labInk, font=\sffamily\small\bfseries, align=left},
    grid style={labGrid, line width=0.3pt},
    major grid style={labGrid, line width=0.3pt},
    axis lines*=left,
    tick align=outside,
    clip=false,
    legend style={draw=none, fill=none, font=\sffamily\footnotesize, text=labSecondary,
      cells={anchor=west}, legend columns=-1, column sep=6pt, at={(0,1.03)}, anchor=south west},
    every axis plot/.append style={line width=1.1pt, line join=round, line cap=round},
    scaled ticks=false,
    /pgf/number format/1000 sep={,},
  },
  lab marks/.style={mark=*, mark size=1.9pt, mark options={draw=white, line width=0.6pt}},
}

\setchessboard{boardfontencoding=LSBC4, boardfontfamily=skaknew, blackfieldcolor=black!18,
  whitefieldcolor=white, setfontcolors, showmover=false, label=true, labelfontsize=6pt,
  boardfontsize=12pt}

% A move or line in UCI, a position in FEN (breakable at a slash), and a run's name.
\newcommand{\uci}[1]{\texttt{#1}}
\newcommand{\fen}[1]{{\footnotesize\path{#1}}}
\newcommand{\run}[1]{\texttt{#1}}
\newcolumntype{L}{>{\raggedright\arraybackslash}X}
\newcolumntype{R}{>{\raggedleft\arraybackslash}X}

\setlength{\parskip}{0.4em}
\setlength{\parindent}{0pt}
\captionsetup{font=small, labelfont=bf}

% from models/<model>/paper/, \input{../preamble} from papers/<name>/.
% Boards come from chessboard and xskak, plots from pgfplots. A paper prints no
% number by hand: `lab paper` writes every number, table and figure under the
% paper's generated/ folder, and paper.tex says what they mean.
\usepackage[LSBC4,T1]{fontenc}
\usepackage[utf8]{inputenc}
\usepackage{lmodern}
\usepackage[margin=1in]{geometry}
\usepackage{microtype}
\usepackage{amsmath}
\usepackage{booktabs}
\usepackage{tabularx}
\usepackage{array}
\usepackage{xcolor}
\usepackage[obeyspaces,spaces]{url}
\usepackage{xskak}
\usepackage{pgfplots}
\usepgfplotslibrary{groupplots}
\usetikzlibrary{positioning,arrows.meta,calc,fit}
\pgfplotsset{compat=1.18}
\usepackage{caption}
\usepackage{subcaption}
\usepackage{enumitem}
\usepackage[section]{placeins}
\usepackage[hidelinks]{hyperref}

% A generated value by its key: generated/values.tex sets them, \V{key} prints
% one, and a key nobody set stops the build.
\makeatletter
\newcommand{\setvalue}[2]{\expandafter\gdef\csname labvalue@#1\endcsname{#2}}
\newcommand{\V}[1]{\@ifundefined{labvalue@#1}%
  {\PackageError{lab}{No generated value named '#1'}{Run `lab paper` for this model again.}}%
  {\csname labvalue@#1\endcsname}}
\makeatother

% Chart colours. A run keeps its colour in every figure of every paper.
\definecolor{labBlue}{HTML}{2A78D6}
\definecolor{labOrange}{HTML}{EB6834}
\definecolor{labAqua}{HTML}{1BAF7A}
\definecolor{labYellow}{HTML}{EDA100}
\definecolor{labMagenta}{HTML}{E87BA4}
\definecolor{labGreen}{HTML}{008300}
\definecolor{labViolet}{HTML}{6250D6}
\definecolor{labGray}{HTML}{898781}
\definecolor{labInk}{HTML}{0B0B0B}
\definecolor{labSecondary}{HTML}{52514E}
\definecolor{labGrid}{HTML}{E1E0D9}
\definecolor{labAxis}{HTML}{C3C2B7}
\definecolor{labWash}{HTML}{CDE2FB}

\pgfplotsset{
  lab/.style={
    font=\sffamily\footnotesize,
    axis line style={labAxis},
    tick style={labAxis},
    tick label style={labSecondary},
    label style={labSecondary},
    title style={labInk, font=\sffamily\small\bfseries, align=left},
    grid style={labGrid, line width=0.3pt},
    major grid style={labGrid, line width=0.3pt},
    axis lines*=left,
    tick align=outside,
    clip=false,
    legend style={draw=none, fill=none, font=\sffamily\footnotesize, text=labSecondary,
      cells={anchor=west}, legend columns=-1, column sep=6pt, at={(0,1.03)}, anchor=south west},
    every axis plot/.append style={line width=1.1pt, line join=round, line cap=round},
    scaled ticks=false,
    /pgf/number format/1000 sep={,},
  },
  lab marks/.style={mark=*, mark size=1.9pt, mark options={draw=white, line width=0.6pt}},
}

\setchessboard{boardfontencoding=LSBC4, boardfontfamily=skaknew, blackfieldcolor=black!18,
  whitefieldcolor=white, setfontcolors, showmover=false, label=true, labelfontsize=6pt,
  boardfontsize=12pt}

% A move or line in UCI, a position in FEN (breakable at a slash), and a run's name.
\newcommand{\uci}[1]{\texttt{#1}}
\newcommand{\fen}[1]{{\footnotesize\path{#1}}}
\newcommand{\run}[1]{\texttt{#1}}
\newcolumntype{L}{>{\raggedright\arraybackslash}X}
\newcolumntype{R}{>{\raggedleft\arraybackslash}X}

\setlength{\parskip}{0.4em}
\setlength{\parindent}{0pt}
\captionsetup{font=small, labelfont=bf}

% Written by `lab paper game-rating`; do not edit by hand.
\looktrue
\chancefalse
\sequencetrue
\contextgbdttrue
\pairstrue
\variantstrue
\shapesfalse
\gamestrue
\featurestrue
\changestrue
\changehandtrue
\changeepochsfalse
\changestartfalse
\changeattacksfalse
\changegeometryfalse
\changeconvfalse
\changehybridfalse
\changelargerfalse
\morefalse
\shownfalse
\shownpairedfalse
\setvalue{band.1.best}{241.0}
\setvalue{band.1.floor}{592.6}
\setvalue{band.1.reference}{352.6}
\setvalue{band.3.best}{181.1}
\setvalue{band.3.floor}{74.7}
\setvalue{band.3.name}{1400-1700}
\setvalue{band.3.reference}{160.1}
\setvalue{band.5.best}{277.1}
\setvalue{band.5.floor}{636.2}
\setvalue{band.5.reference}{391.8}
\setvalue{bands.best-ahead.of}{four of the five bands}
\setvalue{bands.count}{five}
\setvalue{bands.lead.narrowest.name}{1400-1700}
\setvalue{best.beside.closer}{60}
\setvalue{best.beside.correlation}{0.80}
\setvalue{best.close.median}{169.0}
\setvalue{best.close.p90}{422.1}
\setvalue{best.close.within.100}{31}
\setvalue{best.close.within.200}{57}
\setvalue{best.close.within.300}{76}
\setvalue{best.device}{NVIDIA A100-SXM4-80GB}
\setvalue{best.epochs.run}{8}
\setvalue{best.gain}{+38.56}
\setvalue{best.gain.high}{+42.32}
\setvalue{best.gain.low}{+35.06}
\setvalue{best.gain.of-lift}{55}
\setvalue{best.gain.percent}{16}
\setvalue{best.name}{\run{conv-small}}
\setvalue{best.parameters}{1{,}730{,}209}
\setvalue{best.parameters.millions}{1.7}
\setvalue{best.r2}{0.5392}
\setvalue{best.r2.high}{0.5584}
\setvalue{best.r2.low}{0.5182}
\setvalue{best.reliability.groups}{20}
\setvalue{best.reliability.highest.happened}{2058}
\setvalue{best.reliability.highest.said}{2106}
\setvalue{best.reliability.lowest.happened}{960}
\setvalue{best.reliability.lowest.said}{896}
\setvalue{best.reliability.spread}{661}
\setvalue{best.reliability.widest-gap}{64.2}
\setvalue{best.rows-a-second}{2{,}500}
\setvalue{best.score}{204.1}
\setvalue{best.score.high}{208.2}
\setvalue{best.score.low}{200.2}
\setvalue{best.seconds}{5 min}
\setvalue{best.spearman}{0.7324}
\setvalue{best.weights.megabytes}{7.0}
\setvalue{challenger.name}{\run{conv-small}}
\setvalue{code.commit}{\texttt{8c77161}}
\setvalue{context.lift.share}{8}
\setvalue{cost.seconds}{43 min}
\setvalue{cost.seconds.networks}{42 min}
\setvalue{cost.seconds.others}{12 s}
\setvalue{data.fingerprint}{\texttt{e08251266d6e}}
\setvalue{data.games.test}{4{,}999}
\setvalue{data.licence}{CC0 1.0}
\setvalue{data.month}{October 2026}
\setvalue{data.name}{\texttt{games-2026-10}}
\setvalue{data.positions}{2{,}762{,}344}
\setvalue{data.positions.a-row}{27.7}
\setvalue{data.rows}{99{,}699}
\setvalue{data.rows.test}{9{,}974}
\setvalue{data.share.test}{10}
\setvalue{data.share.train}{85}
\setvalue{data.share.valid}{5}
\setvalue{describe.elo-gap.q1}{$-$410}
\setvalue{describe.elo-gap.q25}{$-$39}
\setvalue{describe.elo-gap.q5}{$-$133}
\setvalue{describe.elo-gap.q75}{39}
\setvalue{describe.elo-gap.q95}{133}
\setvalue{describe.elo-gap.q99}{410}
\setvalue{describe.loss-mean.q50}{4.12}
\setvalue{describe.player-elo.q25}{1231}
\setvalue{describe.player-elo.q5}{877}
\setvalue{describe.player-elo.q75}{1792}
\setvalue{describe.player-elo.q95}{2124}
\setvalue{describe.spent-mean.q50}{10.5}
\setvalue{describe.steps.q50}{26}
\setvalue{describe.steps.q95}{48}
\setvalue{examples.count}{three}
\setvalue{floor.constant}{1516.0}
\setvalue{floor.score}{312.3}
\setvalue{grid.loss-pace.conv-small.widest-off}{47}
\setvalue{grid.rating-gap.1400-1700.--200.gbm-context.said}{1536}
\setvalue{grid.rating-gap.1400-1700.--200.gbm-hand.said}{1624}
\setvalue{grid.rating-gap.1400-1700.--200.happened}{1533}
\setvalue{grid.rating-gap.1400-1700.--200.said}{1640}
\setvalue{grid.rating-gap.1400-1700.-100-to-0.said}{1537}
\setvalue{grid.rating-gap.1400-1700.-200-to--100.said}{1600}
\setvalue{grid.rating-gap.1400-1700.0-100.said}{1515}
\setvalue{grid.rating-gap.1400-1700.100-200.said}{1478}
\setvalue{grid.rating-gap.1400-1700.200+.gbm-context.said}{1507}
\setvalue{grid.rating-gap.1400-1700.200+.gbm-hand.said}{1426}
\setvalue{grid.rating-gap.1400-1700.200+.happened}{1552}
\setvalue{grid.rating-gap.1400-1700.200+.said}{1400}
\setvalue{grid.smallest-cell}{30}
\setvalue{label.mean}{1510.5}
\setvalue{label.sd}{382.7}
\setvalue{ladder.at-last-epoch.of}{all three networks}
\setvalue{look.date}{23 November 2026}
\setvalue{look.resamples}{200}
\setvalue{look.smallest-compared}{200}
\setvalue{look.test-rows}{9{,}974}
\setvalue{manifest.games-read}{1{,}689{,}786}
\setvalue{method.gbm.features.percent}{80}
\setvalue{method.gbm.leaves}{127}
\setvalue{method.gbm.patience}{100}
\setvalue{method.gbm.rate}{0.05}
\setvalue{method.gbm.rounds}{4{,}000}
\setvalue{method.gbm.rows.percent}{80}
\setvalue{method.resamples}{1{,}000}
\setvalue{networks.ahead.of}{all three networks}
\setvalue{networks.patience}{three}
\setvalue{networks.warmup.percent}{5}
\setvalue{overall.conv-small.off}{$-$2.5}
\setvalue{overall.gbm-hand.off}{+5.5}
\setvalue{pair.gbm-hand-half.gbm-hand.gain}{$-$1.905}
\setvalue{pair.gbm-hand-half.gbm-hand.gain.high}{$-$0.805}
\setvalue{pair.gbm-hand-half.gbm-hand.gain.low}{$-$2.913}
\setvalue{pair.gbm-hand-half.gbm-hand.gain.percent}{0.79}
\setvalue{pair.gbm-hand-half.gbm-hand.verdict}{significantly worse}
\setvalue{pair.gbm-hand-no-accuracy.gbm-hand.gain}{$-$30.68}
\setvalue{pair.gbm-hand-no-accuracy.gbm-hand.gain.high}{$-$27.15}
\setvalue{pair.gbm-hand-no-accuracy.gbm-hand.gain.low}{$-$34.16}
\setvalue{pair.gbm-hand-no-choices.gbm-hand.gain}{$-$10.82}
\setvalue{pair.gbm-hand-no-choices.gbm-hand.gain.high}{$-$8.76}
\setvalue{pair.gbm-hand-no-choices.gbm-hand.gain.low}{$-$12.66}
\setvalue{pair.gbm-hand-no-tempo.gbm-hand.gain}{$-$6.583}
\setvalue{pair.gbm-hand-no-tempo.gbm-hand.gain.high}{$-$4.791}
\setvalue{pair.gbm-hand-no-tempo.gbm-hand.gain.low}{$-$8.421}
\setvalue{pair.gbm-hand-quarter.gbm-hand.gain}{$-$4.285}
\setvalue{pair.gbm-hand-quarter.gbm-hand.gain.high}{$-$2.883}
\setvalue{pair.gbm-hand-quarter.gbm-hand.gain.low}{$-$5.506}
\setvalue{pair.gbm-hand-quarter.gbm-hand.gain.percent}{1.8}
\setvalue{pair.gbm-hand-quarter.gbm-hand.verdict}{significantly worse}
\setvalue{pair.transformer-medium-hand.transformer-medium.gain}{$-$14.35}
\setvalue{pair.transformer-medium-hand.transformer-medium.gain.high}{$-$11.82}
\setvalue{pair.transformer-medium-hand.transformer-medium.gain.low}{$-$16.72}
\setvalue{pair.transformer-medium-hand.transformer-medium.kept}{not accepted}
\setvalue{pair.transformer-medium-hand.transformer-medium.valid.verdict}{no better}
\setvalue{pair.transformer-medium-hand.transformer-medium.verdict}{significantly worse}
\setvalue{recipe.games}{50{,}000}
\setvalue{recipe.least-moves}{10}
\setvalue{recipe.min-ply}{11}
\setvalue{recipe.most-moves}{48}
\setvalue{reference.close.within.100}{26}
\setvalue{reference.close.within.200}{48}
\setvalue{reference.close.within.300}{67}
\setvalue{reference.features}{20}
\setvalue{reference.leans.1.feature}{mean win probability}
\setvalue{reference.leans.1.share}{16}
\setvalue{reference.leans.2.feature}{mean win probability lost per move}
\setvalue{reference.leans.2.share}{14}
\setvalue{reference.leans.3.feature}{mean number of legal moves}
\setvalue{reference.leans.3.share}{8}
\setvalue{reference.lift.percent}{22}
\setvalue{reference.r2}{0.3741}
\setvalue{reference.reliability.spread}{784}
\setvalue{reference.reliability.widest-gap}{35.1}
\setvalue{reference.score}{242.7}
\setvalue{reference.spearman}{0.5964}
\setvalue{reference.trees}{315}
\setvalue{results.date}{23 November 2026}
\setvalue{run.conv-small.best-epoch}{8}
\setvalue{run.conv-small.epochs}{8}
\setvalue{run.conv-small.parameters.millions}{1.7}
\setvalue{run.conv-small.r2}{0.5392}
\setvalue{run.conv-small.rmse}{259.2}
\setvalue{run.conv-small.score}{204.1}
\setvalue{run.conv-small.valid.score}{206.5}
\setvalue{run.gbm-context.gain}{$-$64.20}
\setvalue{run.gbm-context.r2}{0.0294}
\setvalue{run.gbm-context.score}{306.9}
\setvalue{run.gbm-context.spearman}{0.1607}
\setvalue{run.gbm-hand-no-accuracy.score}{273.4}
\setvalue{run.transformer-medium-hand.best-epoch}{7}
\setvalue{run.transformer-medium-hand.epochs}{8}
\setvalue{run.transformer-medium-hand.gain}{+17.30}
\setvalue{run.transformer-medium-hand.gain.high}{+19.68}
\setvalue{run.transformer-medium-hand.gain.low}{+14.90}
\setvalue{run.transformer-medium-hand.score}{225.4}
\setvalue{run.transformer-medium.gain}{+31.65}
\setvalue{run.transformer-medium.gain.high}{+35.05}
\setvalue{run.transformer-medium.gain.low}{+28.19}
\setvalue{run.transformer-medium.parameters.millions}{11.4}
\setvalue{run.transformer-medium.score}{211.1}
\setvalue{run.transformer-small.gain}{+37.80}
\setvalue{run.transformer-small.gain.high}{+41.24}
\setvalue{run.transformer-small.gain.low}{+34.08}
\setvalue{run.transformer-small.parameters.millions}{2.1}
\setvalue{run.transformer-small.r2}{0.5397}
\setvalue{run.transformer-small.rmse}{259.0}
\setvalue{run.transformer-small.score}{204.9}
\setvalue{run.transformer-small.valid.score}{203.9}
\setvalue{runs.count}{six}
\setvalue{slice.band.-1100.challenger.off}{+220}
\setvalue{slice.band.-1100.reference.off}{+349}
\setvalue{slice.band.1100-1400.challenger.off}{+97}
\setvalue{slice.band.1400-1700.challenger.off}{$-$22}
\setvalue{slice.band.1400-1700.gain}{$-$21.00}
\setvalue{slice.band.1400-1700.gain.high}{$-$13.88}
\setvalue{slice.band.1400-1700.gain.low}{$-$27.89}
\setvalue{slice.band.1700-2000.challenger.off}{$-$127}
\setvalue{slice.band.2000+.challenger.off}{$-$264}
\setvalue{slice.band.2000+.reference.off}{$-$390}
\setvalue{slice.gap.--200.challenger}{270.8}
\setvalue{slice.gap.--200.challenger.off}{+171}
\setvalue{slice.gap.--200.challenger.said}{1479}
\setvalue{slice.gap.--200.happened}{1308}
\setvalue{slice.gap.--200.reference.off}{+200}
\setvalue{slice.gap.--200.reference.said}{1508}
\setvalue{slice.gap.-50-to-0.challenger}{194.8}
\setvalue{slice.gap.0-50.challenger}{193.7}
\setvalue{slice.gap.200+.challenger}{292.8}
\setvalue{slice.gap.200+.challenger.off}{$-$212}
\setvalue{slice.gap.200+.challenger.said}{1527}
\setvalue{slice.gap.200+.gain}{$-$0.50}
\setvalue{slice.gap.200+.gain.high}{+19.81}
\setvalue{slice.gap.200+.gain.low}{$-$19.27}
\setvalue{slice.gap.200+.happened}{1739}
\setvalue{slice.gap.200+.reference.off}{$-$212}
\setvalue{slice.gap.200+.reference.said}{1527}
\setvalue{slice.gap.groups}{eight}
\setvalue{slice.loss.-2.happened}{1576}
\setvalue{slice.loss.12+.happened}{1132}
\setvalue{slice.loss.2-4.happened}{1584}
\setvalue{slice.loss.ahead.of}{all six groups}
\setvalue{slice.moves.-15.challenger}{224.2}
\setvalue{slice.moves.-15.gain}{+25.77}
\setvalue{slice.moves.-15.happened}{1425}
\setvalue{slice.moves.-15.reference}{249.9}
\setvalue{slice.moves.40-48.gain}{+50.48}
\setvalue{slice.moves.48+.challenger}{180.5}
\setvalue{slice.moves.48+.gain}{+50.54}
\setvalue{slice.moves.48+.happened}{1619}
\setvalue{slice.moves.48+.reference}{231.0}
\setvalue{slice.pace.-5.happened}{1486}
\setvalue{slice.pace.30+.happened}{1526}
\setvalue{slice.pace.ahead.of}{all six groups}
\setvalue{slice.result.lost.challenger.off}{$-$14}
\setvalue{slice.result.lost.happened}{1495}
\setvalue{slice.result.won.challenger.off}{+8.8}
\setvalue{slice.result.won.happened}{1504}
\setvalue{slice.side.black.challenger.off}{$-$6.1}
\setvalue{slice.side.white.challenger.off}{+1.0}
\setvalue{slice.start.-600.gain}{+14.05}
\setvalue{slice.start.-600.gain.high}{+37.77}
\setvalue{slice.start.-600.gain.low}{$-$13.05}
\setvalue{slice.start.-600.rows}{220}
\setvalue{slices.ahead}{45}
\setvalue{slices.behind}{one}
\setvalue{slices.compared}{48}
\setvalue{study.features.gbm-hand.least.group}{clock use}
\setvalue{study.features.gbm-hand.most.group}{move accuracy}
\setvalue{study.features.gbm-hand.worse.of}{all three ablations}
\setvalue{task.context.count}{three}
\setvalue{task.hand.count}{17}
\setvalue{task.marks.count}{five}
\setvalue{task.reference}{\run{gbm-hand}}
\setvalue{task.steps.count}{13}


\begin{document}

\labtitle{Game rating}
  {What rating does a player's play in one game look like?}
  {Results of \V{results.date}\quad\textperiodcentered\quad dataset \V{data.name},
   build \V{data.fingerprint}\quad\textperiodcentered\quad code \V{code.commit}}

\begin{summary}
A rating is earned over many games, and one game is a single reading of it.
This paper trains models that take that reading. A model is given a player's
moves through one game, in order: the position before each, both clocks, the
seconds the move took and what an engine says it lost. It says the player's
rating, and reads neither player's rating nor the result. The data is
\V{data.rows} players of \V{recipe.games} rated rapid games on Lichess, split
by game, and the score is the mean absolute error in rating points. The best
of \V{runs.count} runs, \V{best.name}, is off by \V{best.score} points on a
typical game (95\% interval \V{best.score.low} to \V{best.score.high}),
against \V{reference.score} for gradient-boosted trees on
\V{task.hand.count} counts over the player's moves and \V{floor.score} for
one rating said for every player.
\end{summary}

\section{The question}

A rating is a summary of many games. One game is a single reading of it, and
a noisy one: a player rated 1500 has games that look like 1300 and games that
look like 1700. A review of one game can still use the reading. What an
engine says of a game is measured against perfect play, as so many points of
winning chances lost or an accuracy in percent~\cite{lichessaccuracy}, and a
player has no scale for either. A rating is the scale a player already has.
``You played this game like a player two hundred points above you'' is read
at once.

The coaching app this lab works for reviews a player's games one at a time.
A reading of one game in rating points would let a review say what the game
looked like beside the player's usual level: a good day, a bad one, or a
game that was played well on the board and lost on the clock. Before the app
says any of that, three things have to be measured, and this paper is about
them. How far off is one game's reading? What in a game gives the rating
away: what the moves lost, how the clock was used, or the positions
themselves? And whose level is being read? Lichess pairs players of like
rating and every position of a game is made by both of them, so a model can
read the level of the game and never tell its two players apart.

Reading a player's strength off the moves is an old idea. Intrinsic ratings
fit a model of a fallible player to the moves of many games, as an engine
judges them, and place the player on the rating scale~\cite{intrinsic}. A
study of human error against perfect play relates how often a move is a
mistake to the player's rating, the time left and how hard the position
is~\cite{humanerror}. Models of human play run the other way: Maia-2 is told
a rating and says what such a player moves~\cite{maia2}, and Allie reads a
game as a sequence of moves and of the time each took~\cite{allie}. This
paper asks a narrow question with the lab's common tools: one game, one
player, one number.

\begin{table}[htbp]\centering\small
\begin{tabularx}{\textwidth}{lL}
\toprule
Reads & Each of the player's moves from the sixth on, in order, up to
\V{recipe.most-moves} of them: the position before the move, the opponent's
last move, the move played, both clocks, the seconds the move took, the
player's winning chances before it and what it lost, the number of legal
moves, whether it took or gave check, and the castling rights. Beside the
moves, the time control and the game's length. Neither player's rating and
not the result. \\
Predicts & The player's Lichess rapid rating when the game was played. \\
Data & Rated Lichess rapid games of \V{data.month} (\V{data.licence}) with
engine evaluations and clocks. A row is one player in one game, and both
players of a game are rows. \\
Split & By game: \V{data.share.train}\% train, \V{data.share.valid}\% valid,
\V{data.share.test}\% test, so both players of a game are in the same
part. \\
Score & Mean absolute error in rating points on the test rows. Lower is
better. \\
\bottomrule
\end{tabularx}
\caption{The model in brief.}
\label{tab:brief}
\end{table}

\section{The data}

The source is the Lichess database of rated games for
\V{data.month}~\cite{lichessdb}. The builder reads the month as a stream from
its start and stops once \V{recipe.games} games have given rows, so the
sample is the first games of the month and not a draw from all of it. A game
gives rows when it is rated, played at rapid speed between two rated humans,
ended over the board or on time, and carries Lichess's engine analysis and
both clocks.

A row is one player in one game. A move of the player counts from ply
\V{recipe.min-ply} on, which is each side's sixth move, after the opening
moves players know by heart, and it counts when the evaluation before and
after it and the clocks around it are all known. A player with fewer than
\V{recipe.least-moves} such moves gives no row, and a row keeps the first
\V{recipe.most-moves} of them at most. Where both players have enough moves,
both are rows. That gives \V{data.rows} rows that hold \V{data.positions}
positions between them, \V{data.positions.a-row} a row on average
(Table~\ref{tab:data}). The split is a hash of the game, so both players of
a game are on the same side of it. The hash is the one the lab's dataset of
single positions uses, so a game falls in the same part there and here.

\paragraph{The rating.}
The label is the rating the game's record gives the player: their Lichess
rapid rating when the game was played. Lichess keeps its ratings with
Glicko-2~\cite{glicko2}, in which the rating of a new player is still
unsettled. A game's record does not say how settled a rating is, so every
rated player counts alike here.

\paragraph{What a move lost.}
An engine's evaluation is turned into the mover's chance of winning by the
curve Lichess publishes~\cite{lichessaccuracy}, which the app uses as well; a
parity check in this repository holds the lab's curve to the app's. What a
move lost is the mover's chance before it less the chance after it, in points
of a percent. The evaluations are the ones Lichess's own analysis attached to
the game. Lichess runs that analysis when somebody asks for it, so these are
games somebody wanted to look at again.

\paragraph{The seconds a move took.}
Lichess records each player's clock after every move. The time spent on a
move is the mover's clock before it, less the clock after it, plus the
increment the time control adds.

\begin{table}[htbp]\centering\small
\block{data}
\caption{The dataset's rows by split. A row is one player in one game.}
\label{tab:data}
\end{table}

\iflook
The builder read \V{manifest.games-read} games of the month to find the
\V{manifest.games-kept} that gave rows. The ratings run from
\V{describe.player-elo.min} to \V{describe.player-elo.max} around a mean of
\V{describe.player-elo.mean}, and half of them lie between
\V{describe.player-elo.q25} and \V{describe.player-elo.q75}
(Figure~\ref{fig:label}). The two players of
a game are close: the player's rating less the opponent's lies between
\V{describe.elo-gap.q25} and \V{describe.elo-gap.q75} for half of the rows,
and between \V{describe.elo-gap.q5} and \V{describe.elo-gap.q95} for nine in
ten. A row holds \V{describe.steps.q50} of the player's moves at the median,
and one row in four holds \V{describe.steps.q75} or more
(Table~\ref{tab:describe}).

\begin{figure}[htbp]\centering
\block{label}
\caption{The player's rating over every row of the dataset.}
\label{fig:label}
\end{figure}

\begin{table}[htbp]\centering\small
\block{describe}
\caption{The spread of a few columns over every row of the dataset: the mean,
five quantiles, and the least and the most. Times are in seconds, and what
the moves lost is in points of winning chances.}
\label{tab:describe}
\end{table}
\fi

\section{What a model reads}
\label{sec:reads}

% How a position reaches a learner: the same for every model, so each paper
% reads it from here, under its own section heading, and adds what is its own.
Every model of the lab reads a position the same way, so two models differ in
what they are asked and never in how a board reaches them.

\paragraph{The board.}
A position arrives as a FEN and becomes 64 tokens, one a square. The mover's
pawn, knight, bishop, rook, queen and king are tokens 1 to 6, the opponent's
are 7 to 12, and an empty square is 0. When Black is to move the board is
turned and the colours are swapped, so the side to move always plays up the
board. A network then learns one game, not a white one and a black one, and a
position it meets with White to move also teaches it the mirrored position
with Black to move.

\paragraph{Marks.}
What the task knows about particular squares travels as marks: yes-or-no
channels laid over the 64 squares, which a network reads together with the
pieces. This model has \V{task.marks.count} (Table~\ref{tab:marks}). The first
is always the square a pawn may be taken on in passing. The others name the
squares a move leaves and reaches, which is how a network is shown a move
without being told what the move does.

\paragraph{Context numbers.}
Whatever is not about a square is a number beside the board: the four castling
rights, the mover's first, and what the task knows about the game. This model
has \V{task.context.count} (Table~\ref{tab:reads}). Every learner reads them,
the trees that see no board included. A network reads them standardised by the
mean and spread of the train rows.

\paragraph{Hand features.}
Gradient-boosted trees cannot read 64 tokens, so the lab also computes what a
person would count on a board: material, mobility, checks and captures on
offer, pieces that can be taken, passed pawns and pressure on each king, and
for a line of moves what the line does. Table~\ref{tab:hand} in the appendix
lists all \V{task.hand.count}. They serve twice. They are what the reference
run reads, which makes them the bar for a network: a network that reads the
squares themselves should find at least what a few dozen counts already hold.
And a run can hand them to a network beside the board, where they join the
context numbers, which tests whether the network had found them on its own.

\begin{table}[htbp]\centering\small
\block{reads}
\caption{The context numbers: what every learner reads beside the board.}
\label{tab:reads}
\end{table}

\begin{table}[htbp]\centering\small
\block{marks}
\caption{The marks: what the task tells a network about squares. A square
is always given as the mover sees it, after the board is turned.}
\label{tab:marks}
\end{table}

\iflook
\begin{figure}[htbp]\centering
\block{encoding}
\caption{One position of the dataset twice. Left: as it was played, with an
arrow for the opponent's move that led to it (grey) and for the move the row
is about (blue). Right: as a network reads it, turned so that the mover plays
up the board, with a wash on every square a mark is set on: grey for the
opponent's move and, where the task marks a line of moves, blue for the
mover's and orange for the replies.}
\label{fig:encoding}
\end{figure}
\fi


\paragraph{This model reads what an engine said of each move, and neither
rating.}
The marks are the squares of the opponent's last move and of the player's own
move, so a network sees what was played in a position and not only where the
player stood. A move's numbers are of three kinds. The clock: both clocks and
the seconds the move took, as logarithms, and the ply. The engine: the
player's winning chances before the move and what the move lost. The board:
how many legal moves there were, whether the move took or gave check, and the
castling rights. The context numbers are the time control's starting time as
a logarithm, its increment, and the length of the game in plies.

The engine is read on purpose, where the lab's model of think time leaves it
out. That model is a yardstick of what a player can see during a game. This
one is for a review after it, when the analysis is there to be read.

Neither player's rating is read. The pairing would give the answer away,
since the opponent's rating alone lands close to the player's. The result is
not read either: the reading is meant to be of the play.

\section{The learners}
\label{sec:learners}

% The learners every model is given: the same ladder of runs on every task,
% so each paper reads it from here under its own section heading.
A run is one learner, trained on the train rows, tuned on the valid rows and
scored once on the test rows. Every model keeps the same ladder of runs, so
each step up it says what one more thing is worth. Table~\ref{tab:runs} lists
this model's.

\paragraph{The floor.}
\run{constant} says one number for every row: the mean label of the train
rows, or their median where the score is an absolute error. A model that does
not beat it has learned nothing.

\paragraph{Trees.}
\run{gbm-context} is gradient-boosted trees (LightGBM~\cite{lightgbm}) on the
context numbers alone, with no view of the board. It measures how much of the
answer is already in the numbers beside the board. \run{gbm-hand} adds the
hand features and is the \emph{reference}: every other run is compared with
it on the same test rows. Both grow up to \V{method.gbm.rounds} trees of up
to \V{method.gbm.leaves} leaves at a learning rate of \V{method.gbm.rate},
each tree on a random \V{method.gbm.rows.percent}\% of the rows and
\V{method.gbm.features.percent}\% of the features, and stop when the valid
rows have not improved for \V{method.gbm.patience} rounds.

\paragraph{Networks.}
A network is a body, which turns the tokens, the marks and the context numbers
into one vector, and a head of two layers, which turns that vector into the
answer. Two bodies ship, so that a leaderboard can say what the architecture
is worth on a task.

The \emph{square transformer}~\cite{attention} treats each square as a token:
an embedding of its piece, plus a learned embedding of the square, plus a
linear map of the square's marks. Every square attends to every other from the
first layer, so a bishop on a1 and a king on h8 meet at once. Attention is
told about the board's geometry by a learned bias for each offset between two
squares, one per head~\cite{relative}: files, ranks, diagonals and knight hops
are all offsets, so a head can learn to look along them. The context numbers
enter as one more token, and the body's output is that token beside the mean
of the 64 square tokens.

The \emph{convolutional tower} is residual blocks of three-by-three
convolutions over the board, the kind of network the first AlphaZero-style
engines used~\cite{alphazero}. The context numbers are added to every square
after the first layer, and the output is the mean and the maximum over the
squares. It sees further only as it gets deeper. It is cheap and trains on a
CPU, and it is here as the honest alternative: where it matches the
transformer, the cheaper model wins.

Sizes are settings of a run, not code. \run{conv-small} and
\run{transformer-small} are sized to train on a CPU. \run{transformer-medium}
is sized for a GPU, and \run{transformer-medium-hand} is the same network
reading the hand features beside the board. Table~\ref{tab:settings} in the
appendix has every size and schedule.

\paragraph{One training loop.}
Every network is trained the same way: AdamW~\cite{adamw}, a learning rate
that rises over the first \V{networks.warmup.percent}\% of the steps and then
follows a cosine down to zero, gradients clipped, mixed precision on a GPU.
The loss follows the score: squared error for a quantity ranked by RMSE, a
smoothed absolute error for one ranked by MAE, the log loss for a chance. A
quantity's label is standardised by the train rows. After each epoch the
network is scored on the valid rows, the weights of the best epoch are the
ones kept, and training stops early once \V{networks.patience} epochs pass
without a better one.

\begin{table}[htbp]\centering\footnotesize
\block{runs}
\caption{The runs, best first. \emph{Kept} is what the valid rows chose: the
number of trees, or the epoch whose weights were kept. \emph{Time} is the
whole run, training and scoring.}
\label{tab:runs}
\end{table}


\section{How a result is scored}
\label{sec:scoring}

The model's card named its score before the first full run: the mean
absolute error on the test rows, in rating points, where lower is better. A
run beats the reference when its paired gain has a 95\% interval above zero.
Three more scores are reported beside it: the root mean squared error, which
a few large misses move more, R\textsuperscript{2}, the share of the
ratings' variance a run explains, and Spearman's rank correlation, which asks
only whether the players come out in the right order.

Results are also given by the player's rating in \V{bands.count} bands. The
bands are bands of the label itself, so they show whom a run reads well: a
run that pulls every reading toward the middle has an error that grows
toward both ends.

Both players of a game are test rows, and what a model says of one is not
independent of what it says of the other: the two shared every position. The
intervals allow for that.

% How every result of the lab is scored: the same for every model, so each
% paper names its own score under its section heading and reads the rest here.
\paragraph{Noise is measured over games.}
Rows that come from one game are not independent, so every interval in this
paper resamples whole games. The test games are drawn again with replacement
\V{method.resamples} times, the score is computed each time on the rows of the
games drawn, and the interval given is the middle 95\% of those scores.

\paragraph{A gain is paired.}
A run and the reference are scored on the same resampled games, and what is
reported is the difference, with its own interval. A gain counts when that
interval leaves zero out, and is then set in bold. Two runs that are both good
on easy games and bad on hard ones are compared fairly this way, where their
two separate intervals would overlap and hide a real difference.

\paragraph{Test rows are for the final score.}
\ifchance
The number of trees, the epoch whose weights are kept and the calibration of
the chance are all chosen on the valid rows.
\else
The number of trees and the epoch whose weights are kept are both chosen on
the valid rows.
\fi
The test rows are read once a run, for the scores printed here.

\ifchance
\paragraph{A chance is calibrated before it is scored.}
A model can rank well and still say four in ten where three in ten happen. So
one line is put through each run's log-odds, a slope and an intercept fitted
on the valid rows (Platt scaling~\cite{platt}), and applied to its test
predictions before they are scored. A run is then not marked down for a scale
that one line can fix, and what is left of its miscalibration is its own.
\fi

\paragraph{Reading the tables and figures.}
A run keeps one colour and one mark in every figure, and the reference is
always orange. Every figure has its numbers in a table, here or in the
appendix, so nothing has to be read off a colour. No number in this paper was
typed: each is written from the committed record of a run by the command that
builds the paper.


\section{Results}
\label{sec:results}

\subsection{The leaderboard}

Read from the bottom up, Table~\ref{tab:leaderboard} is a ladder. The floor
says \V{floor.constant} for every player, the median rating of the train
rows, and is off by \V{floor.score} points. Trees that read only the time
control and the game's length, \V{run.gbm-context.name}, reach
\V{run.gbm-context.score}, which is \V{context.lift.share}\% of what the
reference takes off the floor. The reference, which also reads the counts
over the player's moves, reaches \V{reference.score}, which is
\V{reference.lift.percent}\% under the floor, and the best run
\V{best.score} (Figure~\ref{fig:scores}).

\begin{table}[htbp]\centering\small
\block{leaderboard}
\caption{The runs on the \V{data.rows.test} test rows, from
\V{data.games.test} games, best first. The gain is paired with
\V{task.reference} on the same resampled games, positive when the run is
better, and bold when its interval leaves zero out.}
\label{tab:leaderboard}
\end{table}

\begin{figure}[htbp]\centering
\block{scores}
\caption{Every run's MAE with its 95\% interval. The vertical line is the
reference's score.}
\label{fig:scores}
\end{figure}

\subsection{What a network adds}

Beyond the noise, \V{ladder.ahead.of} beat the trees and
\V{ladder.behind.of} fell short of them (Figure~\ref{fig:gains}).

\begin{figure}[htbp]\centering
\block{gains}
\caption{Each network's paired gain in MAE over the reference, with its 95\%
interval. Right of the line a network is ahead of the trees, left of it
behind.}
\label{fig:gains}
\end{figure}

\subsection{The order and the number}

The best run's rank correlation with the player's rating is
\V{best.spearman}, and it explains a share of \V{best.r2} of the ratings'
variance (Table~\ref{tab:metrics}).

\begin{table}[htbp]\centering\small
\block{metrics}
\caption{Every score of every run on the test rows. MAE and RMSE are in
rating points and lower is better; for R\textsuperscript{2} and Spearman's
rank correlation higher is better. The floor says one number, so it has no
rank correlation.}
\label{tab:metrics}
\end{table}

\subsection{By the player's rating}

A band rewards hedging. The bands are cut by the player's own rating, so the
floor, which says one rating for every player, has the smallest error of all
in the band that rating lies in and the largest at both ends
(Table~\ref{tab:bands}). A model that leans toward the middle looks good in
the middle for the same reason. So a band's column is read from the top
down, run against run, and never across.

\begin{table}[htbp]\centering\small
\block{bands-table}
\caption{MAE by the player's rating band. Read a column from the top down:
the bands are cut by the truth, so one band cannot be compared with another.}
\label{tab:bands}
\end{table}

\begin{figure}[htbp]\centering
\block{bands}
\caption{Each network's lead over the reference in MAE, band by band: the
differences of the scores in Table~\ref{tab:bands}, without intervals. Above
the dashed line a network is ahead of the trees.}
\label{fig:bands}
\end{figure}

\subsection{How the networks learned}

Figure~\ref{fig:curves} follows each network of the ladder on the valid rows
after every epoch. By that score, \V{ladder.at-last-epoch.of} peaked at the
last epoch given.

\begin{figure}[htbp]\centering
\block{curves}
\caption{Each ladder network's MAE on the valid rows after every epoch. The
dashed line is the reference's score on the same rows.}
\label{fig:curves}
\end{figure}

\subsection{What the trees lean on}

The reference grew \V{reference.trees} trees over \V{reference.features}
features (Figure~\ref{fig:leans}). The three it leans on most are:
\V{reference.leans.1.feature} (\V{reference.leans.1.share}\% of the gain),
\V{reference.leans.2.feature} (\V{reference.leans.2.share}\%) and
\V{reference.leans.3.feature} (\V{reference.leans.3.share}\%). A share of
gain says what the trees used. It does not say what makes a player strong.

\begin{figure}[htbp]\centering
\block{leans}
\caption{The ten features the reference's trees gained most from: each
feature's share of the total gain of the splits made on it.}
\label{fig:leans}
\end{figure}

\ifchanges
\section{One change at a time}
\label{sec:changes}

% One change at a time: the runs that change one thing in another run, set
% against that run, and how to read them. The same for every model, so a
% paper reads this under its own heading, after its results, then says what
% its own numbers mean, then reads studies.tex. Each part waits for its runs:
% a model that has none prints nothing.
\ifpairs
\paragraph{Runs set against the run each changed.}
The ladder says what each learner is worth. A change asks a narrower
question: what is one thing worth inside a learner? Twice the epochs, a body
that starts from another run's weights, the hand features under a network.
Each is a pair of runs that differ in that one thing.
Table~\ref{tab:pairs} sets the first of each pair against the second on the
same resampled test games, the way every run is set against the reference,
and gives the first run's gain. This is where the lab's rule for keeping a
change is read: a change is kept as a run of the model's own when its run is
better on the valid rows and the interval of its gain on the test rows
leaves zero out, however good its own score looks. A change that is not kept
is not thrown away: its row stays in this table, as a measurement of what
the change is worth.

\begin{table}[htbp]\centering\footnotesize
\block{pairs}
\caption{Runs set against the run each changed or has to beat. A gain is
the first run's over the second's, positive when the first is better. On the
valid rows it is the difference of the two scores. On the test rows it is
measured on the same resampled games, and it is bold when its 95\% interval
leaves zero out.}
\label{tab:pairs}
\end{table}
\fi

\ifvariants
Figure~\ref{fig:variants} follows each network of the ladder that another
run changed, beside the runs that changed it, on the valid rows. A run keeps
the colour of the network it changed, and its line is broken by what it
changed: long dashes for a longer schedule, dots for a start from another
run's body.

\begin{figure}[htbp]\centering
\block{variants}
\caption{A network of the ladder (the unbroken line) and the runs that
changed one thing in it (the broken lines), after every epoch on the valid
rows. One panel a network, all on one scale. The dashed level is the
reference's score on the same rows.}
\label{fig:variants}
\end{figure}
\fi


% The two studies every model has: fewer training games, and one group of
% hand features left out at a time. The same for every model, so a paper
% reads this after changes.tex and what it says of its own runs there, and
% adds what its own numbers mean. Each study waits for its runs.
\ifgames
\paragraph{Fewer training games.}
How much of a score is the data, and would more of it help? The games study
trains a run again on half and on a quarter of the training games, and
scores it on the same test rows (Table~\ref{tab:games},
Figure~\ref{fig:games}). The games are chosen by a hash, so the quarter lies
inside the half. A network is given twice the epochs on half the games and
four times on a quarter, so it takes as many steps each time and loses only
the games. The study is read forward: where halving the games costs little,
doubling them would buy little, so the curve says what a larger dataset is
worth before anyone builds one.

\begin{table}[htbp]\centering\footnotesize
\block{study-games}
\caption{The games study. Each run is trained on a share of the training
games and scored on all the test rows. Its gain is paired with the run that
was trained on all the games: below zero the run is worse, and bold when the
95\% interval leaves zero out.}
\label{tab:games}
\end{table}

\begin{figure}[htbp]\centering
\block{study-games-figure}
\caption{The games study as a learning curve: the score on the test rows by
the share of the training games a run was given. Each step to the right
doubles the games.}
\label{fig:games}
\end{figure}
\fi

\iffeatures
\paragraph{A group of hand features left out.}
The hand features come in groups: material, the moves on offer, pieces that
can be taken, passed pawns and attacks near the kings, and for a model that
reads a line of moves, its first move and the line as a whole. The features
study trains a run again without one group at a time
(Table~\ref{tab:features}, Figure~\ref{fig:features}). A run that is worse
without a group was using it, and nothing else the run reads holds the same.
A run that is no worse either made no use of the group or finds the same
elsewhere, in another group or on the board. So the study measures what a
group adds on top of everything else, which is less than what it knows: a
group can matter a great deal on a board and cost nothing here.

\begin{table}[htbp]\centering\footnotesize
\block{study-features}
\caption{The features study. Each run reads every hand feature but one
group. Its gain is paired with the run that reads them all: below zero the
run is worse without the group, and bold when the 95\% interval leaves zero
out. The group whose loss costs most comes first.}
\label{tab:features}
\end{table}

\begin{figure}[htbp]\centering
\block{study-features-figure}
\caption{What leaving a group of hand features out costs the reference: the
paired gain of the run without the group over the run with it, with its 95\%
interval. Left of the line the run is worse without the group.}
\label{fig:features}
\end{figure}
\fi


\ifgames
\paragraph{What more games are worth here.}
Against the trees trained on every training game, the trees on half of them
have a paired gain of \V{pair.gbm-hand-half.gbm-hand.gain}, which is
\V{pair.gbm-hand-half.gbm-hand.verdict}, and the trees on a quarter one of
\V{pair.gbm-hand-quarter.gbm-hand.gain}, which is
\V{pair.gbm-hand-quarter.gbm-hand.verdict} (Table~\ref{tab:games}).
\fi

\iffeatures
\paragraph{Which counts the trees need here.}
In \V{study.features.gbm-hand.worse.of} the trees are worse beyond the noise
without the group. The group that costs most is
\V{study.features.gbm-hand.most.group}, at
\V{study.features.gbm-hand.most.gain}, which is
\V{study.features.gbm-hand.most.verdict}. Without what the moves lost the
paired gain is \V{pair.gbm-hand-no-accuracy.gbm-hand.gain}, without how the
clock was used \V{pair.gbm-hand-no-tempo.gbm-hand.gain}, and without what
there was to choose from \V{pair.gbm-hand-no-choices.gbm-hand.gain}.
\fi
\fi

\section{A closer look}
\label{sec:closer}

% What a closer look is and how it is made: the same for every model. A paper
% reads it under its own section heading, before the cuts that are its own.
\iflook
A leaderboard says which run is best. It does not say where a model is right
and where it is wrong, and one score is a poor guide to whether the app can
use it. So the same test rows are cut several ways, on the machine that has
the dataset and every run's test predictions, and the small tables that come
out are committed beside the leaderboard. Everything in this section and the
next is drawn from them (\V{look.test-rows} test rows, looked at on
\V{look.date}).

\paragraph{Slices.}
The test rows are grouped by one column at a time. For each group a slice
gives how many rows fall in it, what happened there, what a run said on
average, and how the run scored. Two runs are set side by side: \V{challenger.name},
the best run that is not the reference, and the reference. The gain of the
first over the second is measured inside each group as it is for the whole
test set, paired and over resampled games (\V{look.resamples} draws), and is
left out for a group with fewer than \V{look.smallest-compared} test rows.
A slice is where a claim such as ``the network helps most on long lines'' is
checked, not made: a group whose interval holds zero says nothing either way.

\paragraph{Grids.}
Two columns at once: what happened in each cell, and what the run said there,
on one scale of colour so the two can be compared by eye. A cell with fewer
than \V{grid.smallest-cell} test rows is left empty.

\paragraph{Examples.}
Positions are chosen by rule, never by eye: ordinary rows picked by a hash of
their key, the rows the best run is surest of at each end, and its largest
misses each way, \V{examples.count} of each. They show what the numbers are
about. They are not evidence of how often anything happens.
\else
\emph{This build of the paper was drawn before the model's closer look was
run. The closer look cuts the test rows by rating, clock and the shape of a
position, checks what a stated number is worth, and shows example positions.
It is written by \texttt{lab analyse} on the machine that holds the dataset
and every run's test predictions, and this section and the next appear once
its file is committed.}
\fi


\iflook
\subsection{How close a reading lands}

An average error hides how the misses are spread. The best run's median miss
is \V{best.close.median} points, and nine in ten of its readings fall within
\V{best.close.p90}. It is within a hundred points of the player's rating for
\V{best.close.within.100}\% of the test rows, within two hundred for
\V{best.close.within.200}\%, and within three hundred for
\V{best.close.within.300}\% (Table~\ref{tab:closeness}).

\begin{table}[htbp]\centering\small
\block{closeness}
\caption{How close each run's readings land on the test rows: the median
miss, the distance nine in ten readings fall within, and the share within
one, two and three hundred points of the player's rating.}
\label{tab:closeness}
\end{table}

\subsection{What a stated rating is worth}

Figure~\ref{fig:reliability} asks whether a stated rating means what it says.
The test rows are put in \V{best.reliability.groups} equal groups by what the
best run said. The group it rates lowest it calls
\V{best.reliability.lowest.said} on average, and those players' ratings
average \V{best.reliability.lowest.happened}. The group it rates highest it
calls \V{best.reliability.highest.said}, against
\V{best.reliability.highest.happened}. In no group are the two further apart
than \V{best.reliability.widest-gap} points. Around one stated number, the
middle four fifths of the players' ratings span about
\V{best.reliability.spread} points.

\begin{figure}[htbp]\centering
\block{reliability}
\caption{What a run said against the player's rating. The test rows are put
in equal groups by what the run said; each point is one group's average
reading against its players' average rating. The wash is the middle four
fifths of the ratings in each group of the first run. On the diagonal a run
is right on average at that level.}
\label{fig:reliability}
\end{figure}

Turned around, the looseness of one reading is a pull toward the middle.
Figure~\ref{fig:band-bias} groups the test rows by the player's own rating
instead. Players rated \V{band.1.name} are read
\V{slice.band.-1100.challenger.off} points from their rating on average by
\V{challenger.name}, and players rated \V{band.5.name}
\V{slice.band.2000+.challenger.off}. For the trees the two figures are
\V{slice.band.-1100.reference.off} and \V{slice.band.2000+.reference.off}.
Every model with error does this, since what it cannot tell apart it places
near the average. A better model does it less.

\begin{figure}[htbp]\centering
\block{slice-band-bias}
\caption{What the two runs said on average, less the players' average
rating, by the player's rating band. Above the line a run reads a band's
players as stronger than they are, and below it as weaker.}
\label{fig:band-bias}
\end{figure}

\subsection{Whose level is read}
\label{sec:gap}

Lichess pairs players of like rating, and every position of a game is made
by both of them. A model can therefore read the level of the game and say it
for both players. Such a model is right on average and wrong for each player
by half the gap between the two: too low for the stronger player, too high
for the weaker one.

Table~\ref{tab:gap} groups the test rows by the player's rating less the
opponent's. Players rated two hundred points or more above their opponent
are read \V{slice.gap.200+.challenger.off} points from their rating on
average by \V{challenger.name}, and players rated two hundred or more below
theirs \V{slice.gap.--200.challenger.off}. For the trees the two figures are
\V{slice.gap.200+.reference.off} and \V{slice.gap.--200.reference.off}
(Figure~\ref{fig:gap-bias}).

\begin{table}[htbp]\centering\footnotesize
\block{slice-gap}
\caption{The test rows by the player's rating less the opponent's. For each
group: its test rows, the players' average rating, what each of the two runs
said on average and its MAE there, and the first run's paired gain over the
second.}
\label{tab:gap}
\end{table}

\begin{figure}[htbp]\centering
\block{slice-gap-bias}
\caption{What the two runs said on average, less the players' average
rating, by the player's rating less the opponent's. Above the line a run
reads a group's players as stronger than they are, and below it as weaker.}
\label{fig:gap-bias}
\end{figure}

That cut alone does not settle whose level is read, because the pull toward
the middle leans the same way. A player far above the opponent is more often
a strong player, and a pull toward the middle reads a strong player too low
whoever the opponent is. Figure~\ref{fig:rating-gap} holds the two apart.
Each row of it is one band of the player's own rating, and along a row only
the opponent changes. A run that reads the player and nothing else says the
same along a row. A run that reads the level of the game says less from left
to right, as the opponent gets weaker. The opposite lean is possible too: a
stronger opponent sets harder problems, and a player who loses more to them
looks weaker than against an easier one.

\begin{figure}[htbp]\centering
\block{grid-rating-gap}

\medskip
\block{grid-rating-gap-said}
\caption{The players' average rating, by the player's rating band and the
player's rating less the opponent's. Above: the ratings themselves. Below:
what the best run that is not the reference said. Both grids share one scale
of colour, and a cell with too few test rows is left empty.}
\label{fig:rating-gap}
\end{figure}

\subsection{By what the moves lost and how long they took}

Table~\ref{tab:loss} groups the test rows by what the player's moves lost on
average. Players whose moves lost under two points of winning chances a move
are rated \V{slice.loss.-2.happened} on average, and players whose moves
lost twelve or more \V{slice.loss.12+.happened}. By the seconds a move took
on average, the players' ratings run from \V{slice.pace.-5.happened} for
under five seconds a move to \V{slice.pace.30+.happened} for thirty or more
(Table~\ref{tab:pace}).

\begin{table}[htbp]\centering\footnotesize
\block{slice-loss}
\caption{The test rows by the points of winning chances the player's moves
lost on average.}
\label{tab:loss}
\end{table}

\begin{table}[htbp]\centering\footnotesize
\block{slice-pace}
\caption{The test rows by the seconds the player's moves took on average.}
\label{tab:pace}
\end{table}

Figure~\ref{fig:grid} crosses the two. The upper grid is the players'
average rating and the lower one what the run said. If the lower grid looks
like the upper one, the run has the joint pattern and not only each margin.

\begin{figure}[htbp]\centering
\block{grid-loss-pace}

\medskip
\block{grid-loss-pace-said}
\caption{The players' average rating, by the points of winning chances their
moves lost on average and the seconds a move took on average. Above: the
ratings themselves. Below: what the best run that is not the reference said.
Both grids share one scale of colour, and a cell with too few test rows is
left empty.}
\label{fig:grid}
\end{figure}

\subsection{By the game}

A row can hold as few as \V{recipe.least-moves} of the player's moves and as
many as \V{recipe.most-moves}, and a reading has more to go on in a long
game. Table~\ref{tab:moves} groups the test rows by how many moves were
read. The other cuts of the game are in the appendix: how it ended for the
player (Table~\ref{tab:slice-result}), the player's colour
(Table~\ref{tab:slice-side}), the time control
(Tables~\ref{tab:slice-start} and~\ref{tab:slice-increment}) and what the
player's costliest move lost (Table~\ref{tab:slice-worst}).

\begin{table}[htbp]\centering\footnotesize
\block{slice-moves}
\caption{The test rows by how many of the player's moves the row holds.}
\label{tab:moves}
\end{table}

\subsection{Beside the reference}

Table~\ref{tab:beside} sets each run's readings beside the reference's on
the same test rows: their correlation, and how often the run is the closer
of the two to the player's rating.

\begin{table}[htbp]\centering\small
\block{beside}
\caption{Each run's readings beside the reference's on the same test rows:
their correlation, and how often the run is the closer of the two to the
player's rating.}
\label{tab:beside}
\end{table}

\section{Games}
\label{sec:positions}

The rows below were chosen by rule from the test games. A row is a whole
game of one player, and one board stands for it: the position before the
player's costliest move, seen from the player's side. The grey arrow is the
opponent's last move and the blue arrow is the move the player made. Beside
each board are those two moves, the player's rating, what the best run and
the reference said, and what the game came with. The games are public
Lichess games; no game is named.

\examples{examples-typical}{Ordinary rows: the first \V{examples.count} test
rows in the order of a hash of their key, so nothing about them picked
them.}{fig:typical}

\examples{examples-over}{The best run's largest misses upward: the test rows
where it said the highest rating against the player's own.}{fig:over}

\examples{examples-under}{The best run's largest misses downward: the test
rows where it said the lowest rating against the player's own.}{fig:under}
\fi

\section{What the model does not say yet}
\label{sec:limits}

\begin{itemize}
\item \textbf{One game is a noisy reading.} A player rated 1500 has games
that look like 1300 and games that look like 1700. The error left after the
best run is partly the model's and partly the game's, and nothing here
separates the two. Several games of one player would.
\item \textbf{The pairing shows through.} A position is made by both
players, and Section~\ref{sec:gap} measures how much of a reading is the
level of the game. The same model asked for the opponent's rating would say
it from the other side, and that run is not made yet.
\item \textbf{One number for the game.} A player can open like a master and
finish like a beginner. A reading for the opening, the middlegame and the
ending needs a label for each part, and a rating is a label for none of them.
\item \textbf{The engine is read.} What a move lost comes from the
evaluations Lichess attached. A rating read from the moves and the clock
alone is a different run.
\item \textbf{Which games.} Only analysed games give a row, and players
choose which games to analyse. The evaluation is as deep as Lichess's own
analysis goes. A long game is cut at the player's \V{recipe.most-moves}th
counted move, and the opening is not read at all.
\item \textbf{Whose scale.} The ratings are Lichess rapid's. A reading is on
that scale, which is not the scale of the player's own site, and a fit of
one to the other is a private check.
\end{itemize}

\section{Cost and how to draw this again}
\label{sec:cost}

The \V{runs.count} runs took \V{cost.seconds} in all: \V{cost.seconds.networks}
for the networks, which trained on a rented GPU, and \V{cost.seconds.others}
for the floor and the trees. Table~\ref{tab:runs} has every run's time and
Table~\ref{tab:settings} the networks' settings.

Everything is drawn from committed files. The dataset is built by
\texttt{lab data games} from the recipe, the runs are trained by
\texttt{lab all game-rating}, each run writes its record under the model's
\texttt{results/} folder, and \texttt{lab paper game-rating} writes every
number, table and figure of this paper from those records and compiles it. A
record names the build of the data it read (\V{data.fingerprint}) and the
commit of the code that trained it (\V{code.commit}). Weights and test
predictions are not committed. They are kept on the owner's machine.

\section{What comes next}

\begin{itemize}
\item \textbf{The opponent's rating as the label.} The same rows and the
same learners, asked for the other player: how well a run does there is how
much of a reading is the game's and not the player's.
\item \textbf{Several games of one player.} One game is a noisy reading,
and a review would read a player's last twenty. How fast the error falls
with the number of games needs rows that say which games one player played,
which this dataset does not keep.
\item \textbf{A reading for each part of the game.} The idea behind this
model is a rating for the opening, the middlegame and the ending.
\item \textbf{Without the engine.} The moves and the clocks alone, to say
how much of the reading the engine's judgement is.
\item \textbf{A body that has read boards before.} In the lab's models of
one position, a network that starts from a body trained on engine
evaluations is ahead of the same network started from nothing. A network
over a game cannot start from one yet.
\item \textbf{The player's own games.} A fit of this model's readings to the
player's rating on the player's own site, kept private, before the app
states a reading.
\end{itemize}

\bibliographystyle{unsrt}
\bibliography{../../../papers/references}

\clearpage
\appendix

\section{Sizes, schedules and features}

\begin{table}[htbp]\centering\footnotesize
\block{settings}
\caption{The ladder networks' sizes and schedules, from their run configs.
\emph{Layers over} and \emph{heads over} are those of the transformer that
reads a row's positions in order, above the body that reads each one.
\emph{Rate} is the peak learning rate. A setting they all share is in the
line under the table.}
\label{tab:settings}
\end{table}

\begin{table}[htbp]\centering\footnotesize
\block{hand}
\caption{The hand features: counts over the player's moves, which
\V{task.reference} reads beside the context numbers, and a network when its
run asks for them.}
\label{tab:hand}
\end{table}

\iflook
\section{The extremes of what the model says}

\examples{examples-highest}{The test rows for which the best run said the
highest rating.}{fig:highest}

\examples{examples-lowest}{The test rows for which the best run said the
lowest rating.}{fig:lowest}

\section{Every cut of the test rows}

Each table below groups the test rows by one column. The columns are the
group, its test rows, the players' average rating, then for each of the two
runs the rating it said on average and its MAE, and the paired gain of the
first run over the second with its 95\% interval, bold when the interval
leaves zero out.

\block{slices}
\fi

\end{document}
